\chapter*{Abstract}
\addcontentsline{toc}{chapter}{Abstract}

The transition from informal mathematical language to machine-verifiable logic represents a new frontier in formal verification, known as Auto-Formalization. This thesis, inspired by Google DeepMind's AlphaProof and Harmonic's Aristotle, explores the development of a "Truth Scanner"—a neuro-symbolic pipeline designed to translate English mathematical definitions and theorems into Lean 4 code. 

The project follows a structured 9-week roadmap. Initial phases focused on mastering the Lean 4 tactic language through the completion of the Natural Number Game. The core of the research involves selecting key definitions from standard undergraduate texts, such as Rosen’s \textit{Discrete Mathematics}, and mapping them to formal logical components. 

A primary contribution of this work is the implementation of a "Repair Loop" that utilizes Large Language Models (LLMs) to perform self-correction by feeding Lean compiler errors back into the translation prompt. By evaluating the pipeline's ability to resolve ambiguities and logical gaps often found in human-written textbooks, this thesis aims to create a "Wikipedia of Truth" where mathematical claims are backed by rigorous, machine-checked proofs.

\textbf{Keywords:} Lean 4, Auto-Formalization, LLMs, Formal Verification, Neuro-Symbolic Pipelines, Discrete Mathematics.